% ------------------------------------------------------------------------------
% INTRODUCCIÓN
% ------------------------------------------------------------------------------
\section{Introducción}

Durante la presente unidad se abordó el estudio y uso de \LaTeX, un sistema de composición de textos de alta calidad tipográfica orientado especialmente a la creación de documentos técnicos y científicos. Su utilidad principal radica en la automatización del formato, la gestión impecable de referencias bibliográficas y la escritura avanzada de expresiones matemáticas, alejándose del paradigma WYSIWYG (\textit{What You See Is What You Get}) propio de los procesadores de texto tradicionales, para acercarse a un enfoque estructural basado en código.

A nivel personal, el lenguaje \LaTeX{} ya formaba parte de mis herramientas de uso cotidiano. Durante mi educación media, desarrollé de manera independiente un apunte íntegro de Probabilidad y Estadística de más de 60 páginas, lo que me otorgó un dominio sólido sobre entornos matemáticos, teoremas, y control tipográfico. Sin embargo, esta unidad resultó fundamental para adquirir una competencia que hasta ahora no había explorado: \textbf{la creación de presentaciones mediante la clase \texttt{beamer}}.

Aprender a trasladar la lógica estática de un texto continuo a un formato dinámico de diapositivas interactivas, manejando retículas espaciales, entornos como \texttt{columns} y controlando el flujo visual de la exposición oral, constituyó el mayor valor agregado y el principal aprendizaje de este módulo formativo.

\newpage

% ------------------------------------------------------------------------------
% TRABAJOS REALIZADOS
% ------------------------------------------------------------------------------
\section{Trabajos Realizados}

La unidad estuvo compuesta por tres laboratorios de complejidad incremental, orientados a dominar distintas facetas de la escritura académica en ingeniería:

\begin{itemize}
	\item \textbf{Laboratorio 1 (Formato Matemático):} El objetivo inicial fue estructurar documentos base utilizando jerarquías de secciones y comandos de formato. Se trabajó en la maquetación de un instrumento de evaluación (Control 2 de Matemáticas). Este requirió el uso intensivo de entornos anidados (\texttt{enumerate}, \texttt{itemize}) y modo matemático avanzado (\texttt{amsmath}) para representar lugares geométricos, parábolas e hipérbolas.

	\item \textbf{Laboratorio 2 (Leyes de Newton y Bibliografías):} Orientado a la redacción de un artículo estructurado (\texttt{article}), el objetivo fue incorporar elementos de síntesis y gestión documental. Se redactó un informe sobre las Leyes de Newton, lo que implicó la creación de tablas dimensionalmente ajustadas (\texttt{tabular}) para relacionar conceptos físicos y su aplicación en programación. Adicionalmente, se implementó un sistema de referencias automatizado, vinculando el documento a un archivo externo \texttt{.bib} mediante el gestor \texttt{BibTeX}.

	\item \textbf{Laboratorio 3 (Presentaciones en Beamer):} El desafío final consistió en utilizar la clase \texttt{beamer} para diseñar un material de apoyo visual. Se elaboró una presentación sobre la biografía de Isaac Newton y sus postulados. La maquetación requirió distribuir el contenido en múltiples columnas (\texttt{\textbackslash begin\{columns\}}), insertar recursos gráficos (\texttt{\textbackslash includegraphics}) y aplicar una jerarquía visual estricta mediante el uso de bloques semánticos (\texttt{block}) para separar ecuaciones, descripciones y ejemplos de la vida cotidiana.
\end{itemize}

\subsection{Metodología y Flujo de Trabajo Híbrido}

A diferencia del flujo de trabajo estándar exclusivamente en la nube, diseñé una metodología híbrida adaptada a la naturaleza de cada evaluación. Para los \textbf{trabajos individuales} (pre y post laboratorios), utilicé un entorno local sobre el sistema operativo Arch Linux. El código fue redactado en el editor modal \texttt{Neovim}, potenciado con \texttt{Vimtex} y \texttt{LuaSnip} (para la expansión rápida de macros), compilado en tiempo real con \texttt{latexmk} y renderizado mediante el visor de documentos \texttt{Zathura}. Posteriormente, el código era versionado en \texttt{GitHub} y, finalmente, subido a \texttt{Overleaf} de forma empaquetada (\texttt{.zip}) para cumplir con los lineamientos institucionales de entrega.

Por el contrario, los \textbf{laboratorios grupales} fueron desarrollados de forma íntegra en la plataforma \texttt{Overleaf} para favorecer el trabajo colaborativo síncrono. En esta dinámica establecimos roles claros: mi compañero se enfocó en la investigación y redacción de contenidos, mientras que yo asumí el rol de maquetador, encargándome de la estructura del documento, inserción de imágenes y gestión de referencias. Una vez finalizados, los proyectos eran descargados, recompilados en mi entorno local para control de calidad, y respaldados en mi repositorio de \texttt{GitHub}.

\subsection{Evidencia de los laboratorios}

A continuación, se adjunta evidencia del trabajo realizado, destacando el entorno de desarrollo local empleado luego de la sincronización con la plataforma institucional.

% Asegúrate de tener estas imágenes en tu carpeta img/
\insertimage[\label{img:evidencia1}]{neovim.png}{width=0.8\textwidth}{Entorno de trabajo local (Neovim + Zathura) durante la compilación del código fuente del Laboratorio 2.}

\insertimage[\label{img:evidencia2}]{nvim.png}{width=0.8\textwidth}{Entorno de trabajo local (Neovim + Zathura) para la presentación diseñada con la clase Beamer (Laboratorio 3).}

\newpage
% ------------------------------------------------------------------------------
% PRINCIPALES DIFICULTADES
% ------------------------------------------------------------------------------
\section{Principales Dificultades}

A lo largo del desarrollo de las actividades, las dificultades enfrentadas no radicaron en la sintaxis del lenguaje, sino en aspectos de gestión de tiempo, colaboración y dependencias locales:

\begin{enumerate}
	\item \textbf{Gestión del tiempo (Laboratorio 2):} La extensión de la actividad fue considerable. El proceso de investigar el material en Internet, resumirlo respetando la estructura pedida, formatear las secciones y las imágenes, y gestionar la correcta vinculación de las referencias en el archivo \texttt{.bib}, demandó más tiempo del disponible en el bloque de clases. Esto me obligó a finalizar la actividad de manera asíncrona en mi domicilio.

	\item \textbf{Adaptación del flujo de trabajo colaborativo:} Al realizar los laboratorios de forma grupal, se presentó el desafío de integrar metodologías dispares. Si bien mi entorno óptimo involucra el control de versiones local (\texttt{Git}) y \texttt{Neovim}, mi compañero no estaba familiarizado con estas herramientas. Para evitar fricciones, fue necesario migrar mi flujo hacia la plataforma \texttt{Overleaf}, asumiendo una pérdida temporal en la velocidad de escritura en pos del trabajo colaborativo.

	\item \textbf{Gestión de Paquetes Locales (Troubleshooting):} Al compilar este mismo informe final utilizando plantillas avanzadas en mi máquina local, me enfrenté a un error de compilación relacionado con una secuencia de control no definida (\texttt{\textbackslash AtBeginShipout}). Tras realizar un proceso de \textit{debugging}, identifiqué que el error derivaba de la ausencia de carga del paquete \texttt{atbegshi} en mi distribución local de \TeX Live. Este problema técnico fue resuelto forzando la importación del paquete en el preámbulo del archivo principal, lo que demuestra que trabajar en local requiere una administración proactiva de dependencias de la cual Overleaf exime al usuario.
\end{enumerate}

\newpage

% ------------------------------------------------------------------------------
% CONCLUSIONES
% ------------------------------------------------------------------------------
\section{Conclusiones}

La finalización de esta unidad ha permitido consolidar y formalizar mis conocimientos previos sobre \LaTeX{}, demostrando empíricamente por qué este sistema de preparación de documentos es el estándar ineludible en el ámbito de las ciencias y la ingeniería. Más allá de la pulcritud estética, el paradigma de separar el contenido de su representación visual permite al estudiante enfocarse plenamente en la lógica y veracidad de la información, delegando al compilador la responsabilidad tipográfica.

De manera particular, el descubrimiento y asimilación de la clase \texttt{beamer} representa un salto cualitativo en mis capacidades de comunicación académica. Entender que un mismo entorno de código puede generar tanto un \textit{paper} formal estructurado mediante referencias \texttt{BibTeX}, como un soporte visual dinámico para una exposición oral, entrega una versatilidad profesional invaluable. Estas herramientas serán fundamentales para el futuro de la carrera, asegurando que la transmisión del conocimiento mantenga siempre el rigor y la claridad que exige la Facultad.
